\section{Constant kernels on five points}
\label{sec:five-point}

The reflected inequality extends beyond the support-size boundary of the
covering bound when the kernel is constant.

\begin{theorem}\label{thm:five-point}
Let $F\subset\mathbb Z$ have five elements. For every nonzero finitely
supported $f\ge0$ and every nonzero $g\ge0$ supported on $F$,
\[
 \|f\star\widetilde g\|_1^4
 < |4F|\,\|f\star g\|_1^4.
\]
In particular, the five-point set in Proposition~\ref{b-five-point}
satisfies this inequality despite failing the sufficient covering bound.
\end{theorem}

We first prove a bound that applies to every five-point support, then
separate Sidon supports from those with pair-sum collisions.

\subsection{A bound from energy and the range of the weights}

\begin{lemma}\label{five-ratio}
For every five-element $F\subset\mathbb Z$ and every $f,g$ as in
Theorem~\ref{thm:five-point},
\[
 \frac{\|f\star\widetilde g\|_1}{\|f\star g\|_1}\le\frac{62}{25}.
\]
Consequently, Theorem~\ref{thm:five-point} holds whenever $|4F|\ge38$.
\end{lemma}
\begin{proof}
For a nonempty finite set $A$, let $r_+$ and $r_-$ count the representations
of a point as $a+p$ and $a-p$, respectively, with $(a,p)\in A\times F$.
Put $P=|A+F|$, $Q=|A-F|$ and $W=5|A|$. The two quadratic energies are
equal: exchanging the two $F$ coordinates bijects the solutions of
$a+p=a'+p'$ with those of $a-p=a'-p'$. Write their common value as $E$.
Since $1\le r_-(s)\le5$ on its support, $r_-(s)^2\le6r_-(s)-5$.
Cauchy--Schwarz applied to $r_+$ gives
\[
 \frac{W^2}{P}\le E\le6W-5Q,
 \qquad
 \frac QP\le\frac{6t-t^2}{5}\le\frac95,
 \quad t=\frac WP.
\]
Integrate this inequality over the level sets of $f$ to obtain
\begin{equation}\label{five-energy}
 \|f\star1_{-F}\|_1\le\frac95\|f\star1_F\|_1.
\end{equation}
For $g>0$ on $F$, put $r=\max_Fg/\min_Fg\ge1$. Pointwise comparison
with the indicator kernels gives
\begin{equation}\label{five-range-energy}
 \frac{\|f\star\widetilde g\|_1}{\|f\star g\|_1}\le\frac95r.
\end{equation}

The dual in Lemma~\ref{lem:cover}, with $k=1_F$, has objective
$S=\sum_iw_i$, where $w_i=g_i z_i$ and $g_i=g(p_i)>0$ for
$F=\{p_1,\ldots,p_5\}$. Its constraints imply
\begin{equation}\label{five-dual-pairs}
 0\le w_i\le1,\qquad
 \frac{g_j}{g_i}w_i+\frac{g_i}{g_j}w_j\le1,\qquad
 w_iw_j\le\frac14\quad(i\ne j).
\end{equation}
The middle inequality retains two nonnegative terms of the constraint
at $p_i+p_j$; AM--GM gives the last one.

Every four coordinates have sum at most two. If their largest value
$u$ is at most $1/2$, this is immediate. Otherwise their sum is at most
$u+3/(4u)\le2$ for $1/2\le u\le1$. It follows that $w_i\ge S-2$
for every $i$. Choose distinct indices at which $g$ takes its maximum
and minimum; when $r=1$, any two distinct indices suffice. Their
constraint in \eqref{five-dual-pairs} gives
\[
 1\ge(r+r^{-1})(S-2),\qquad
 S\le2+\frac1{r+r^{-1}}.
\]
Together with the covering lemma, this proves
\begin{equation}\label{five-range-cover}
 \frac{\|f\star\widetilde g\|_1}{\|f\star g\|_1}
 \le\min\left\{\frac95r,\ 2+\frac1{r+r^{-1}}\right\}.
\end{equation}
For $1\le r\le4/3$, the first bound is at most $12/5$.
For $r\ge4/3$, the second is at most $2+12/25=62/25$.

For nonnegative nonzero $g$, apply the bound to $g+\epsilon1_F$ and
let $\epsilon\downarrow0$. For fixed $f$, both max-convolutions are
supported in fixed finite sets and their values are continuous in $\epsilon$.
The limiting denominator is positive. Finally,
\[
 \left(\frac{62}{25}\right)^4
 =\frac{14776336}{390625}<38,
\]
which proves the last assertion.
\end{proof}

\subsection{Supports with a pair-sum collision}

For an ordinary-autoconvolution-feasible $x\ge0$, write $w_p=x(p)^2$.
The universal constraints are $0\le w_p\le1$ and
$w_pw_q\le1/4$ for distinct points. We use repeatedly the elementary
inequality
\begin{equation}\label{five-box}
 b+c\le a+\frac{bc}{a}\qquad(0\le b,c\le a, a>0),
\end{equation}
which follows from $(a-b)(a-c)\ge0$.

\begin{lemma}\label{five-parallelogram}
If a five-element integer set $F$ contains four distinct points forming a
parallelogram, then $N_F(1_F)\le2$.
\end{lemma}
\begin{proof}
Label the four points so that $p_0+p_3=p_1+p_2$, and label the fifth
point $p_4$. The extra convolution constraint is
\[
 \sqrt{w_0w_3}+\sqrt{w_1w_2}\le\frac12.
\]
If all weights are at most $1/2$, apply \eqref{five-box} to the two
opposite pairs. With $p=\sqrt{w_0w_3}$ and $q=\sqrt{w_1w_2}$,
\[
 w_0+w_1+w_2+w_3\le1+2(p^2+q^2)
 \le1+2(p+q)^2\le\frac32.
\]
Adding $w_4\le1/2$ proves the bound in this region.

Otherwise let $u\in[1/2,1]$ be a largest weight, and put
$a=1/(4u)\in[1/4,1/2]$. Every other weight is at most $a$.
If $u=w_4$, then
\[
 \sum_iw_i\le u+2a+\frac{p^2+q^2}{a},
 \qquad 0\le p,q\le a,\quad p+q\le\frac12.
\]
On this polygon the maximum of $p^2+q^2$ is
$a^2+(1/2-a)^2$, as follows by checking its vertices. Hence
\[
 \sum_iw_i\le2u+\frac1u-1\le2.
\]
The last function is convex and equals two at both endpoints.

If the largest weight is a corner, relabel so that $u=w_0$ and put
$p=\sqrt{w_1w_2}\in[0,a]$. The opposite weight is at most
$(1/2-p)^2/u$. Equation~\eqref{five-box} and $w_4\le a$ give
\[
 \sum_iw_i\le u+2a+\frac{p^2}{a}+\frac{(1/2-p)^2}{u}.
\]
This is convex in $p$, so the endpoint bounds are
\[
 A(u)=u+\frac{3}{4u},\qquad
 B(u)=u+\frac1u-\frac1{4u^2}+\frac1{16u^3}.
\]
Both are convex on $[1/2,1]$; for the second,
$B''(u)=(2u^2-3u/2+3/4)/u^5>0$.
Their endpoint values are $(2,7/4)$ and $(2,29/16)$, respectively.
Thus every feasible profile has total at most two.
\end{proof}

\begin{lemma}\label{five-progression}
If a five-element integer set $F$ contains a three-term progression,
then $N_F(1_F)\le17/8$.
\end{lemma}
\begin{proof}
Write $b,c$ for the endpoint weights and $v$ for the middle weight.
The extra constraint is
\begin{equation}\label{five-ap-constraint}
 v+2\sqrt{bc}\le1.
\end{equation}
If every weight is at most $1/2$, then \eqref{five-box} gives
\[
 b+c+v\le\frac12+2bc+v
 \le\frac12+\frac{(1-v)^2}{2}+v
 =1+\frac{v^2}{2}\le\frac98.
\]
The other two weights contribute at most one.

Otherwise let $u\in[1/2,1]$ be a largest weight and set $a=1/(4u)$.
All remaining weights are at most $a$. There are three cases.

\paragraph{An endpoint is largest.}
Take $b=u$. For $v\in[0,a]$, the other endpoint is at most
$(1-v)^2/(4u)\le a$, so
\[
 \sum w\le u+2a+v+\frac{(1-v)^2}{4u}.
\]
Convexity in $v$ reduces this to $v=0,a$, giving $A(u)$ above and
\[
 C(u)=u+\frac1u-\frac1{8u^2}+\frac1{64u^3}.
\]
The function $C$ is convex because
$C''(u)=(2u^2-3u/4+3/16)/u^5>0$.
Its endpoint values are $17/8$ and $121/64$; also $A(u)\le2$.

\paragraph{The middle point is largest.}
Take $v=u$. Equations~\eqref{five-box} and \eqref{five-ap-constraint}
give $b+c\le a+u(1-u)^2$. Thus
\[
 \sum w\le D(u)=u+\frac{3}{4u}+u(1-u)^2.
\]
Here $D''(u)=6u-4+3/(2u^3)>0$ on $[1/2,1]$, and its endpoint
values are $17/8$ and $7/4$.

\paragraph{An outside point is largest.}
The other outside weight is at most $a$, and
\[
 b+c+v\le v+a+\min\left\{a,\frac{(1-v)^2}{4a}\right\},
 \qquad 0\le v\le a.
\]
If $u\in[3/4,1]$, the simple bound $b+c+v\le3a$ gives
$\sum w\le u+1/u\le25/12<17/8$.
If $u\in[1/2,3/4]$, then $a\in[1/3,1/2]$. The two branches
meet at $v_0=1-2a\in[0,a]$. The first branch increases to one.
The second is convex, with endpoint values one and
$2a+(1-a)^2/(4a)$. Therefore
\[
 \sum w\le\max\left\{u+\frac1{4u}+1,\
                       2u+\frac{13}{16u}-\frac12\right\}.
\]
Both functions are convex on $[1/2,3/4]$. Their endpoint values are
$(2,25/12)$ and $(17/8,25/12)$, respectively. This completes the cases.
\end{proof}

\begin{proposition}\label{five-nonsidon}
If a five-element integer set $F$ is not Sidon, then
$N_F(1_F)^4<|4F|$.
\end{proposition}
\begin{proof}
If $F$ contains a four-point parallelogram, Lemma~\ref{five-parallelogram}
gives $N_F(1_F)^4\le16<17\le|4F|$. The last bound follows from
repeated application of $|A+B|\ge|A|+|B|-1$ in $\mathbb Z$.

Otherwise every pair-sum collision consists of one diagonal pair and one
off-diagonal pair, forming a three-term progression. A sum cannot have two
different off-diagonal pairs, since they would be disjoint and form a
parallelogram. Each point is the centre of at most one progression, and
the two extreme points are never centres. Thus there are at most three
collisions among the fifteen unordered pair sums, giving
\[
 |2F|\ge12,\qquad |4F|\ge2|2F|-1\ge23.
\]
Lemma~\ref{five-progression} now gives
$N_F(1_F)^4\le83521/4096<23$.
\end{proof}

\subsection{Five-point Sidon supports}

\begin{proposition}\label{five-sidon-size}
Every five-element Sidon subset $F$ of $\mathbb Z$ satisfies $|4F|\ge39$.
Equality is attained by $F=\{0,2,7,8,11\}$.
\end{proposition}
\begin{proof}
We use Freiman's classical $3k-4$ theorem in the following form:
if $A\subset\mathbb Z$, $|A|=k\ge3$ and $|A+A|\le3k-4$, then $A$
is contained in an arithmetic progression with at most
$|A+A|-k+1$ elements. This is the equal-set specialization of the
statement recalled by Bollob\'as, Leader and Tiba~\cite[p.~1]{blt}.

Translate $F$ and divide by the greatest common divisor of its
differences, writing
\[
 F=\{0,a,b,c,d\},\qquad 0<a<b<c<d,\qquad \gcd(F)=1.
\]
These operations preserve the Sidon property and sumset cardinalities.
Set $B=F+F$. Then $|B|=15$, $\min B=0$, $\max B=2d$, and
$\gcd(B)=1$ because $F\subset B$.
If $|4F|=|B+B|\le38$, the theorem applies, since $38\le3\cdot15-4$.
It puts $B$ in an arithmetic progression with at most $38-15+1=24$
elements. Its step divides every difference of $B$, so its absolute
value is one. Consequently $2d\le23$, and $d\le11$.

The ten positive differences of a five-point Sidon set are distinct,
so $d\ge10$. If $d=10$, those differences would be $1,\ldots,10$,
whose sum is odd. Their sum is instead $4d+2c-2a$, which is even.
Thus the supposition $|4F|\le38$ forces $d=11$.

We classify this last case by hand. The ten differences omit one number
$m$ from $1,\ldots,11$, and the parity argument shows that $m$ is even.
Put $P_i=\{i,11-i\}$ for $1\le i\le5$. The six endpoint differences
\[
 a,b,c,11-a,11-b,11-c
\]
occupy exactly three distinct pairs $P_i$. The three remaining
differences $b-a,c-b,c-a$ form a Schur triple: the sum of the smaller
two is the largest. They must be three of the four members of the two
unused pairs, say $P_r,P_s$ with $r<s$.

Deleting one element of the ordered list $r,s,11-s,11-r$ must therefore
leave a Schur triple. Deleting the smallest cannot work, since
$s+(11-s)=11>11-r$. The other deletions give $s=2r$, $2r+s=11$, or
$r+2s=11$. Their complete solution list is
\begin{center}
\begin{tabular}{cccc}
\toprule
Omitted $m$ & Unused $(r,s)$ & Interior differences & $c-a$\\
\midrule
$2$  & $(1,2)$ & $\{1,9,10\}$ & $10$\\
$4$  & $(2,4)$ & $\{2,7,9\}$  & $9$\\
$6$  & $(3,5)$ & $\{3,5,8\}$  & $8$\\
$8$  & $(3,4)$ & $\{3,4,7\}$  & $7$\\
$10$ & $(1,5)$ & $\{1,5,6\}$  & $6$\\
\bottomrule
\end{tabular}
\end{center}
The interior points must occupy the other three complementary pairs,
once each. For $m=2$, the span $c-a=10$ is impossible inside $[1,10]$.
For $m=4$, it forces $(a,c)=(1,10)$, using the same complementary pair.
For $m=6$, the possibilities $(a,c)=(1,9),(2,10)$ give only
$b=4,7$, respectively. For $m=8$, the possible endpoint pairs are
$(1,8),(2,9),(3,10)$; the outer two use the forbidden $P_3$, and the
middle one repeats $P_2$. For $m=10$, the choices $a=1,4$ use the
forbidden $P_1$. The remaining choices $(a,c)=(2,8),(3,9)$ give only
$b=7,4$, respectively. Up to reflection, the only sets are therefore
\[
 F_1=\{0,1,4,9,11\},\qquad F_2=\{0,2,7,8,11\}.
\]

For the first representative,
\[
 2F_1=\{0,1,2,4,5,8,9,10,11,12,13,15,18,20,22\}.
\]
Successive elements are at most three apart, so
$2F_1+\{0,1,2\}$ contains $[0,24]$. The rest of $[0,38]$ is supplied by
\[
\begin{aligned}
\relax[25,28]&=15+[10,13],& [29,33]&=20+[9,13],\\
 [34,35]&=22+[12,13],& 36&=18+18,\\
 37&=15+22,&38&=18+20.
\end{aligned}
\]
All summands belong to $2F_1$. Also $40=18+22$, $42=20+22$ and
$44=22+22$. The only other candidates are $39,41,43$; each requires a
summand $20$ or $22$, but its complementary summand is absent. Hence
\[
 4F_1=[0,38]\cup\{40,42,44\},\qquad |4F_1|=42.
\]
For the second representative, Proposition~\ref{b-five-point} gives
\[
 4F_2=\{0,2,4\}\cup[6,38]\cup\{40,41,44\},\qquad |4F_2|=39.
\]
Here all intervals are integer intervals. Both possibilities contradict
$|4F|\le38$, and the second proves sharpness.
\end{proof}

The proof classifies normalized span $11$ only. It does not classify all
equality cases for $|4F|=39$.


\begin{proof}[Proof of Theorem~\ref{thm:five-point}]
For non-Sidon supports, combine Proposition~\ref{five-nonsidon} with
Lemma~\ref{lem:cover}. For Sidon supports, use
Proposition~\ref{five-sidon-size} and Lemma~\ref{five-ratio}.
\end{proof}

For $f=1_A$ and $g=1_F$, Theorem~\ref{thm:five-point} gives
$|A-F|^4<|A+F|^4|4F|$ for every finite nonempty integer set $A$ and every
five-element integer set $F$. The theorem also permits arbitrary
nonnegative weights in $f$ and $g$; the kernel $k$ is fixed to $1_F$.
